\documentclass[UTF8, xcolor=table, compress]{beamer}
%\usepackage{fontspec}
%\setsansfont{宋体}
% \setCJKmainfont[BoldFont={Adobe Heiti Std},ItalicFont={Adobe Kaiti Std}]{Adobe Song Std}
% \setCJKmainfont[BoldFont={Source Han Sans},ItalicFont={Adobe Kaiti Std}]{Source Han Serif}

\usepackage{latexsym,amssymb,amsmath,amsbsy,amsopn,amstext,xcolor,multicol,mathtools}
\usepackage{graphicx,wrapfig,fancybox}
\usepackage{pgf}
\usepackage{zjubeamer}
%\usepackage[backend=bibtex,style=IEEE,sorting=none]{biblatex} % [参考文献格式](https://www.sharelatex.com/blog/2013/07/31/getting-started-with-biblatex.html)
\usepackage[backend=bibtex,sorting=none]{biblatex} % [参考文献格式](https://www.sharelatex.com/blog/2013/07/31/getting-started-with-biblatex.html) %mac IEEE not found
\usepackage{array}
\usepackage{bm}
\usepackage{caption}
\usepackage[caption=false,font=scriptsize]{subfig}
\usepackage{multirow}
\usepackage{booktabs}
\usepackage{tikz}
\usetikzlibrary{decorations.pathreplacing}
\usepackage{tikzscale}
\usepackage{animate}
%\usepackage{times} %与上面的冲突，加上这个 粗体斜体就失效
%\usepackage{mathptmx}


\defbibheading{bibliography}[\bibname]{} %avoid printbibliography 自动生成目录
\addbibresource{ref/papers-bib-in-google.bib}
\addbibresource{ref/chinese-ref.bib}
%\setbeamertemplate{bibliography item}{\insertbiblabel} %将列表中默认的丑陋的icon 改成数字,或者下面这个也行
\setbeamertemplate{bibliography item}[text] % [ref](http://tex.stackexchange.com/questions/68080/beamer-bibliography-icon)
%\setbeamertemplate{footline}[frame number]{}

%\setframeofframes{of}

\usepackage{boxedminipage} %for: bvh border
\def\fourgraphicswidth{0.35} %0.3\textwidth

\usepackage[ruled]{algorithm2e} %%format of the algorithm
% \usepackage{algpseudocode}
% \floatname{algorithm}{算法}
% \renewcommand{\algorithmicrequire}{\textbf{输入:}} %%Use Input in the format of Algorithm
% \renewcommand{\algorithmicensure}{\textbf{输出:}} %%UseOutput in the format of Algorithm
%\algrenewcommand{\algorithmiccomment}[1]{\hskip3em $\rightarrow$ #1}
% \algrenewcommand{\algorithmiccomment}[1]{ $//$ #1}

\usepackage{listings}
\renewcommand\lstlistingname{代码}
\renewcommand\lstlistlistingname{代码}

\lstset{framexleftmargin=1.4em,
        xleftmargin=1.8em,
        basicstyle=\ttfamily\small,
        %frame=shadowbox, numberstyle=\tiny, breaklines=true,
        frame=single,
        numberstyle=\tiny, breaklines=true,
        keywordstyle=\color{blue!70}\bfseries,
        %commentstyle=\color{red!50!green!50!blue!50},
        rulesepcolor=\color{red!20!green!20!blue!20},
        numbers=none,fontadjust=true}
\lstdefinelanguage{shader}{morekeywords={uniform, layout, uniform, vec2, vec3, vec4, in, out, gl_Position, dot, flat, int ,float, gl_VertexID, xyz, w, x, y, z, location, version, sampler2DRect, bgr, gl_FragData, texture2DRect, gl_TexCoord,for,xy},morecomment=[l]{//}}

%\setbeameroption{show notes} %un-comment to see the notes

%\usepackage{pgfpages}
%\renewcommand\pgfsetupphysicalpagesizes{%
%    \pdfpagewidth\pgfphysicalwidth\pdfpageheight\pgfphysicalheight%
%}
%\setbeameroption{show notes on second screen}

\begin{document}

\setbeamerfont{footnote}{size=\tiny}
\setbeamerfont{caption}{size=\scriptsize}
\setbeamertemplate{caption}[numbered]
\setbeamerfont{subsection in toc}{size=\footnotesize}
\renewcommand*{\bibfont}{\footnotesize}

\graphicspath{{figures/}}

\title{对内存友好的确定性增量梯度方法}
\author[谢嘉豪]{(申请浙江大学工学学士学位论文答辩报告)\vskip 20pt答辩人：谢嘉豪\vskip 5pt指导教师：钱~徽~教授}
\institute[浙江大学~计算机科学与技术学院]{\small \vskip 70pt}
%\date{2015-06-07}
\date{\small \vskip -17pt二〇一八年六月}
%\date{\today}



\frame{
\vspace{-15mm}
\titlepage
\vspace{-43mm}
\begin{figure}[htbp]
  \begin{center}
	\includegraphics[width=0.16\linewidth]{ZJU_Logo}
  \end{center}
\end{figure}
%\beign{picture}(1,1)
%\put(6,8){\includegraphics[width=0.15\linewidth]{Tsinghua_University_Logo.eps}}
%\end{picture}
}

\section{}  % a empty section to hide from the navigation bar
\frame {
  \frametitle{目录}
 % \begin{multicols}{2}
  \tableofcontents
%\end{multicols}
  \note{
    我将按照下面如下的次序来介绍本人的工作，首先是课题背景和相关工作，然后介绍凸包围体生成算法和基于凸包围体的碰撞检测算法，最后进行总结。
  }
}

\AtBeginSubsection[] {
\frame<handout:0> {
\frametitle{目录}
\tableofcontents[current,currentsubsection]
  }
  \addtocounter{framenumber}{-1}  %目录页不计算页码
}

\section{介绍}
\miniframesoff

%\subsection{课题背景和意义}
\frame
{
\frametitle{课题背景和意义}
\begin{block}{增量梯度方法}
  \begin{itemize}
    \item 机器学习领域中应用最广泛的一类优化方法
    \item 每次迭代只需计算少量的数据样本所对应的目标函数梯度
    \item 优点：迭代速度快、计算和存储复杂度较低、扩展性好等等
    \item 按抽取数据样本的方式可分为两类：
      \begin{itemize}
        \item 随机方法——随机梯度下降（SGD）及其变种
        \item \textbf{确定性方法}
      \end{itemize}
  \end{itemize}
\end{block}

\pause

\begin{block}{选题意义}
  \begin{itemize}
    \item 理论意义：有助于理解一般增量梯度方法的收敛性质
    \item 实践意义：有助于提高机器学习算法的实用性和可扩展性
  \end{itemize}
\end{block}
}

\frame
{
\frametitle{研究现状}
\footnotesize
问题：$\underset{w \in \mathbb{R}^d}{\mathrm{min}} \ F(w) \coloneqq \frac{1}{N} \sum_{i=1}^N f_i(w) $， \pause 增量梯度方法：$w_{t+1} = w_t - \eta_t \ g(w_t, i_t)$ \\~\\
\pause 其中$g$：真实梯度$\nabla F(w_t)$的估计量，$i_t$：第$t$次迭代选择的数据样本

\pause

\begin{block}{}
选择数据样本：
\begin{itemize}
    \item 随机采样
    \begin{itemize}
        \item \alert<6>{可放回采样}：均匀采样、重要性采样、自适应采样
        \item 不可放回采样：随机重排
    \end{itemize}
    \item 确定性顺序
\end{itemize}
\end{block}

\pause

\begin{block}{}
设计梯度估计方式：
\begin{itemize}
    \item 单独近似：$g(w_t, i_t) = \nabla f_{i_t}(w_t)$
    \item \alert<6>{平均近似}：$g(w_t, i_t) = \frac{1}{N}\sum_{i=1}^N\nabla f_i(w_{\tau_i})$
    % \tikz[remember picture] \node[coordinate,yshift=0.5em] (n1) {};
    \item \alert<6>{混合近似}：$g(w_t, i_t) = \nabla f_{i_t}(w_t) - \nabla f_{i_t}(\tilde{w}) + \nabla F(\tilde{w})$
    \item \alert<6>{……}
    % \tikz[remember picture] \node[coordinate] (n2) {};
\end{itemize}
\end{block}

% \only<5>{\alert{
% \begin{tikzpicture}[overlay,remember picture]
%     \path (n2) -| node[coordinate] (n3) {} (n1);
%     \draw[thick,decorate,decoration={brace,amplitude=3pt}]
%           (n1) -- (n3) node[midway, right=4pt] {降方差方法};
% \end{tikzpicture}}}
}

\frame
{
\frametitle{本文工作}
\begin{itemize}
  \item 提出了“基于快照的增量梯度方法”（SIG）及其加速变种 \vspace*{1em}
  \item 证明了所提出的SIG方法达到线性收敛 \vspace*{1em}
  \item 将SIG算法推广到更一般的单调算子零点问题上，证明了在这一类问题上同样达到线性收敛
\end{itemize}
}

\miniframeson

\section{本文提出的算法}
% \miniframesoff

\subsection{问题和假设}
\frame
{
\frametitle{问题和假设}
\begin{block}{问题}
有限和形式的凸优化问题：
$\underset{w \in \mathbb{R}^d}{\text{min}}\ F(w) = \frac{1}{N}\sum_{i=1}^{N} f_i(w)$
\end{block}

\pause

\begin{block}{假设}
\begin{itemize}
    \item \textbf{（Lipschitz光滑假设）}\
    $\forall i \in \{1, \ldots, N\}, \forall w_1, w_2 \in \mathbb{R}^d$,
    $$
    \| \nabla f_i(w_1) - \nabla f_i(w_2) \| \le L \| w_1 - w_2 \|.
    $$

    \item \textbf{（强凸假设）}\
	$F(w)$是$\mu$-强凸函数，即对任意的$w_1, w_2 \in \mathbb{R}^d$，
    $$
	    (\nabla F(w_1) - \nabla F(w_2))^T(w_1 - w_2) \ge \mu \| w_1 - w_2 \|^2.
    $$

    \item 定义问题的条件数为$\kappa = L / \mu$

\end{itemize}
\end{block}
}

\frame
{
\frametitle{应用例子}
\begin{itemize}
    \item 经验风险最小化问题：带$l_2$惩罚项的最小二乘、逻辑回归、光滑SVM……
    $$
        F(w) = \frac{1}{N}\sum_{i=1}^N\phi(w^\top x_i, y_i) + \frac{\mu}{2}\|w\|^2
    $$
    \item 非凸优化的重要子问题：
    $$
	\min_{w} \frac{1}{N}\sum_{i=1}^{N}\phi_i(w) + \frac{L}{2}\|w - \tilde{w}\|^2
    $$

\end{itemize}
}

\subsection{SIG算法}
\frame
{
\frametitle{\subsecname~ }
基于快照的确定性增量梯度（Snapshot-based Incremental Gradient，简称 SIG）方法：

{
\setlength{\interspacetitleruled}{-.4pt}%
\begin{algorithm}[H]
\SetKwInOut{Input}{input}
\SetKwInOut{Return}{return}
\SetAlgoLined
\Input{$\eta$, $w_0^0$, S} \
\For{$s = 0, 1, \ldots, S-1$}
{
    \alert<2>{compute $G = \frac{1}{N} \sum_{i=1}^{N} \nabla f_i(w_0^s)$} \

    \For{$t = 0, \ldots, N-1$}
    {
    \alert<3>{$i_t = t+1$} \

    %$i_t = \sigma^s(i+1)$, $\sigma^s$ can be any peretatation on $\{1,\ldots,N\}$\

    $w_{t+1}^{s} = w_t^s - \eta \big(\alert<4>{\nabla f_{i_t}(w_t^s) - \nabla f_{i_t}(w_0^s) + G} \big)$ \

    }
    $w_0^{s+1} = w_N^s$ \
}
\Return{$w_0^S$}
\end{algorithm}
}

}

\frame
{
\footnotesize
\frametitle{\subsecname~ }
\begin{block}{SIG算法的特点}
    \begin{itemize}
        \item 以确定性的顺序访问数据
        \item 不需要维护历史的梯度信息
        \item 在光滑和强凸假设下达到线性收敛
    \end{itemize}
\end{block}

\pause

\begin{block}{与SVRG\uncover<2->{\footnote{Johnson \& Zhang, Accelerating stochastic gradient descent using predictive variance reduction. NIPS 2013.}}算法的关系}
    \begin{itemize}
        \item 使用几乎相同的梯度估计量
        \item SVRG：随机访问样本数据，SIG：循环顺序访问
        \item SVRG：内循环长度依赖于问题的条件数，SIG：把样本数作为内循环长度
    \end{itemize}
\end{block}
}

\frame
{
\frametitle{与现有方法的比较}
\footnotesize
\begin{table}[H]
	\centering
	\caption{SGD、SVRG、 SAGA、 DIG、 IAG和SIG方法的比较}
	\label{table:methods}
	\begin{tabular}{|c|c|c|c|c|c|>{\columncolor[gray]{0.8}}c|}
		\hline
		算法 & SGD & SVRG & SAGA & DIG & IAG & SIG \\ \hline
		\begin{tabular}[c]{@{}c@{}}数据访问方式\\  (读/写)\end{tabular}              & 只读   & 只读    & 读+写  & 只读  & 读+写 & 只读   \\ \hline
		\begin{tabular}[c]{@{}c@{}}组成函数选择方式 \\ (随机/确定性)\end{tabular} & 随机   & 确定性    & 随机性    & 确定性  & 确定性   & 确定性   \\ \hline
		存储复杂度 & $\mathcal{O}(d)$ & $\mathcal{O}(d)$ & $\mathcal{O}(Nd)$ & $\mathcal{O}(d)$ & $\mathcal{O}(Nd)$ & $\mathcal{O}(d)$ \\ \hline
		线性收敛？ & $\times$  & \checkmark  &  \checkmark  & $\times$ & \checkmark & \checkmark \\ \hline
	\end{tabular}
\end{table}

}

\subsection{SIG-M算法}
\frame
{
% \footnotesize
SIG-M：引入动量（Momentum）\footnote{Allen-Zhu, Katyusha: The first direct acceleration of stochastic gradient methods. ACM STOC, 2017.}来加速SIG的收敛

\pause

\begin{algorithm}[H]
  \SetKwInOut{Input}{input}
  \SetKwInOut{Return}{return}
  \SetAlgoLined
  \Input{$\alpha$, $w_0^0$, $S$, $\tau_1$, $\tau_2$} \
  $y_0^0 = z_0^0 \leftarrow w_0^0$; \

  \For{$s = 0, 1, \ldots, S-1$}
  {
   compute $G = \frac{1}{N} \sum_{i=1}^{N} \nabla f_i(w_0^s)$ \

   \For{$t = 0, \ldots, N-1$}
   {
    $i_t = t+1$ \

    %$i_t = \sigma^s(i+1)$, $\sigma^s$ can be any peretatation on $\{1,\ldots,N\}$\

    \alert<4>{$w_{t+1}^{s} = \tau_1 z_t^s + \tau_2 w_0^s + (1 - \tau_1 - \tau_2) y_t^s$} \

    \alert<3>{$z_{t+1}^s = z_t^s - \alpha \big(\nabla f_{i_t}(w_t^s) - \nabla f_{i_t}(w_0^s) + G \big)$} \

    \alert<3>{$y_{t+1}^s = w_{t+1}^s + \tau_1 (z_{t+1}^s - z_t^s)$} \
    % $y_{t+1}^s = w_{t+1}^s - \frac{1}{3 L} \big(\nabla f_{i_t}(w_t^s) - \nabla f_{i_t}(w_0^s) + G \big) $ \

   }
   $w_0^{s+1} = w_N^s$, $y_0^{s+1} = y_N^s$, $z_0^{s+1} = z_N^s$ \

  }
  \Return{$w_0^S$}
\end{algorithm}
}



% \miniframeson

\section{理论结果}
\subsection{线性收敛结果}
  \frame
  {
    \frametitle{扰动分析框架}
    \footnotesize
    % SIG内循环迭代：    $w_{t+1}^{s} = w_t^s - \eta \big(\nabla f_{i_t}(w_t^s) - \nabla f_{i_t}(w_0^s) + \frac{1}{N} \sum_{i=1}^{N} \nabla f_i(w_0^s) \big)$ \\~\\

% \pause

    由SIG内循环迭代公式和循环顺序的数据访问可得% ，则对$t = 0, \ldots, N-1$求和得：
    $$
    \begin{aligned}
    w_0^{s+1} &= w_0^s - \eta\sum_{t=0}^{N-1} \nabla f_{t+1} (\alert<4->{w_t^s}) &\;\; \text{SIG} \\
\pause
    w_0^{s+1} &= w_0^s - \eta\sum_{i=1}^{N} \nabla f_{i} (\alert<4->{w_0^s}) &\;\; \text{梯度下降}
    \end{aligned}
    $$ \\~\\

  \uncover<5>
  {
    于是，可以将SIG当做\alert{受扰动的梯度下降}来分析
  }

  }

  \frame
  {
    \begin{block}{}
    \frametitle{重要引理：梯度误差上界}
    在光滑假设下，令SIG算法的步长$\eta$满足$\eta < \frac{1}{NL}$，那么
    $$
      \sum_{t=0}^{N-1} \| \nabla f_{i_t}(w_t^s)\! -\! \nabla f_{i_t}(w_0^s) \|
      \le \frac{\eta L^2 N^2}{2(1 - \eta L N)} \| w_0^s\! -\! w^* \|
    $$ \\~\\
    \end{block}

    \pause
    % \footnotesize
    这说明梯度估计量受到的扰动将随算法收敛而逐渐趋于0
  }


  \frame
  {
    \frametitle{定理1}
    \begin{block}{}
    在强凸和光滑假设下，令步长$\eta \le \mu/(2L^2N)$，则迭代参数到最优点的距离$\|w_0^s - w^*\|$线性收敛：
    $$
    \| w_0^{s+1} - w^* \| \le \alpha \| w_0^{s} - w^* \|,
    $$
    其中$\alpha = 1 - \frac{1}{2} \eta \mu N < 1$.
    \end{block}

\pause

    \footnotesize
    \begin{itemize}
        \item 定理1不仅适用于确定性的样本选择顺序，也适用于任意排列顺序
        % \item SIG算法达到$\epsilon$精度的最坏复杂度（迭代次数）为$\mathcal{O}(N\kappa^2 \text{log}(1/\epsilon))$， \\
        % 而IAG算法的最坏复杂度为$\mathcal{O}(N^2 \kappa \text{log}(1/\epsilon))$
    \end{itemize}
  }

\subsection{推广：单调算子}
  \frame
  {
    \frametitle{单调算子的零点问题}
    \footnotesize
    SIG算法可推广到求解单调算子的零点问题：$0 \in B(w) = \frac{1}{n}\sum_{i=1}^n B_i(w)$ \\~\\

    \pause

    \textbf{假设：}
    \begin{itemize}
        \item 假设每个$B_i(\cdot)$均有Lipschitz常数$L > 0$，即
        $$
        \| B_i(w_1) - B_i(w_2) \|_2 \le L \| w_1 - w_2 \|_2 \;\; \forall w_1, w_2 \in \mathbb{R}^d
        $$
        \item $B(\cdot)$是$\mu$-强单调算子，即
        $$
        \langle B(w_1) - B(w_2), w_1 - w_2 \rangle \ge \mu \| w_1 - w_2 \|^2, \;\; \forall w_1, w_2 \in \mathbb{R}^d
        $$
    \end{itemize}

    \pause
    \begin{block}{}
    \textbf{应用：}鞍点问题，即$\underset{w \in \mathbb{R}^d}{\text{min}} \ \underset{u \in \mathbb{R}^p}{\text{max}} F(w, u)$
    \begin{itemize}
        % \item $F$对第一个变量是凸函数，对第二个变量是凹函数
        \item 具体应用：二元分类模型的AUC问题，图像处理中的去噪、去抖动问题
    \end{itemize}
    \end{block}
  }

  \frame
  {
    \frametitle{定理2}
    \begin{block}{}
    在Lipschitz和强单调假设下，令步长$\eta \le \mu/(2L^2N)$，则迭代参数到最优点的距离$\|w_0^s - w^*\|$线性收敛：
    $$
    \| w_0^{s+1} - w^* \| \le \beta \| w_0^{s} - w^* \|,
    $$
    其中$\beta = 1 - \frac{1}{4} \eta \mu N < 1$.
    \end{block}
  }

\section{数值实验}
\miniframesoff
  \frame
  {
    \frametitle{\secname~ }

  }
\miniframeson

\section{总结与展望}
\miniframesoff
  \frame
  {
    \frametitle{\secname~ }
    \footnotesize
    \begin{block}{总结}
    \begin{itemize}
        \item 提出了确定性增量梯度方法SIG和SIG-M
        \item 证明了SIG达到线性收敛，并将其推广到单调算子的零点问题
        \item 实验结果表明，SIG在大规模优化问题上对内存友好，优于现有的算法；
              动量的引入对SIG算法起到了加速的作用
    \end{itemize}
    \end{block}

    \pause

    \begin{block}{未来工作}
    \begin{itemize}
        \item 改进SIG的理论收敛率到$\mathcal{O}(N\kappa \mathrm{log}(1/\epsilon))$
        \item SIG-M的收敛性分析
        \item 将SIG推广到更一般的带非光滑正则项的优化问题
        \item 设计和分析异步或分布式的变种算法
    \end{itemize}
    \end{block}
  }

  % \frame
  % {
  %   \frametitle{致谢}
  %   \footnotesize
  %   \begin{itemize}
  %       \item 感谢钱徽教授对本文选题和写作的指导
  %       \item 感谢沈泽邦学长、张超学长、周腾飞学长对本文提出的宝贵意见以及在凸优化理论方面对我的指导
  %       \item 感谢实验室的其他同学对我工作的帮助
  %       \item 感谢我的家人对我学业的支持
  %   \end{itemize}
  % }

  \frame
  {
    \frametitle{Q \& A}
    \begin{block}{Questions?}
      ~\\ ~\\
      \center{\Large{Thank you!}}
      \\ ~\\ ~\\ ~\\ ~\\
    \end{block}
  }


\end{document}
